\subsection{Why exceptions, duties and permissions need different logics}

A threshold calculation asks whether a proposition follows from evidence. A duty
asks what an actor ought to do in those circumstances. Treating both as ordinary
implication loses the possibility of noncompliance: an instruction to conduct an
annual review is not evidence that the review occurred. Input/output logic makes
this separation explicit. A normative system transforms factual conditions into
required outputs, without requiring every input to be an output or permitting
unrestricted contraposition \citep[secs.~2--3]{inputoutput2000}.

Let $A$ contain factual propositions and let $G$ contain pairs $(a,x)$ expressing
that condition $a$ calls for outcome $x$. Writing $\operatorname{Cn}$ for classical
consequence, the paper's simplest operation is
\begin{equation}
 \operatorname{out}_1(G,A)
 =\operatorname{Cn}\!\left(G\!\left(\operatorname{Cn}(A)\right)\right).
 \label{eq:io}
\end{equation}
First derive the input conditions, then apply the normative pairs, then unpack
their consequences. Section 3.2 characterizes this operation by strengthening
inputs, conjoining outputs and weakening outputs. A pair connecting ``review
due'' with ``review performed'' supplies a required outcome; it does not create
a recorded fact of performance. Other variants change disjunction and reuse.
LegalMath borrows this separation rather than implementing all input/output logics.

Exceptions ask a different question: which applicable rule defeats another?
\citet[sec.~2]{defeasible2000} distinguish facts, strict rules, defeasible rules,
defeaters and an explicit superiority relation. Their proof theory records both
definite and defeasible provability, and explicit failure to establish either.
A defeasible conclusion needs supporting premises, failure of definite contrary
proof and a response to every applicable opposing rule. Different supporting
rules may defeat different opponents. This ``team defeat'' differs materially
from selecting one table row or finding one Catala exception.

An unretrieved agreement is therefore not a defeater merely because the database
has no row. The product must distinguish explicit negative evidence, a
closed-world assumption and a reviewed default. LegalMath 0.1 chooses nested,
source-linked defaults with conflicts between competing sibling exceptions. It
does not implement the full defeasible proof theory. General priorities remain
an extension requiring a stated defeat policy and conformance tests.

Permissions add another distinction. A weak permission is obtained when an
opposite obligation cannot be established; a strong permission has explicit
normative support. \citet[secs.~2--3, 6]{permissions2012} distinguish permissive
rules from defeaters and distinguish obligation from factual fulfillment.
Their algorithms transform a finite theory while preserving its defeasible
extension. The soundness, completeness and linear-complexity results in section 5
concern that theory and specified data-structure assumptions, not arbitrary
regulations. For the bank, failure to derive a prohibition must never by itself
become a positive trade authorization.

A remedy is different again. The logic of violations of
\citet[secs.~2--5]{violations2006} uses an ordered reparation connective: perform
the primary duty; if it is violated, a further duty can become relevant.
Replacing that structure with ordinary disjunction can make the remedy appear
an equally acceptable way to satisfy the original requirement. A warning delivered
after a deadline can be useful remediation without making the original control
timely. LegalMath retains the breach assessment and records subsequent performance
separately; general compensation and sanction reasoning remain outside its first
event profile.
